\documentclass[../DoAn.tex]{subfiles}
\begin{document}
\begin{center}\bfseries\large TÓM TẮT NỘI DUNG BÀI TẬP LỚN\end{center}
\vspace{0.5cm}
Đề tài xuất phát từ nhu cầu giảng viên tìm và đăng ký phòng học cho hoạt động giảng dạy, họp và công việc chuyên môn. Khi thông tin phòng và lịch sử dụng được quản lý bằng file Excel, việc tìm phòng trống, kiểm tra sức chứa và thiết bị, cũng như đăng ký phòng còn phụ thuộc nhiều vào thao tác thủ công. Nhóm phát triển RoomHub để số hóa quy trình tìm phòng, đặt phòng, phê duyệt, thông báo và cấp quyền truy cập cho giảng viên và cán bộ quản trị. Trong phạm vi pre-MVP, hệ thống được mô hình hóa và thử nghiệm trên một tòa nhà tám tầng có cách tổ chức phòng tương tự tòa C7, Đại học Bách khoa Hà Nội; chưa triển khai trên toàn khuôn viên trường.

RoomHub gồm ứng dụng Android React Native/TypeScript, backend Java 17/Spring Boot cung cấp REST API và cơ sở dữ liệu H2 lưu phía máy chủ qua JPA. Cấu hình hiện tại bật chế độ remote: app trao đổi dữ liệu với backend qua API; chế độ local dùng Async Storage vẫn được giữ riêng. Đây là hai chế độ dữ liệu, không phải đồng bộ nền hai chiều giữa bản local offline và server. Các phép kiểm tra quy tắc được tổ chức ở tầng dịch vụ; thao tác ghi phía máy chủ chạy trong giao dịch có khóa bi quan.

Với phòng gắn khóa số, quản trị viên hoặc người dùng đã được cấp quyền theo phòng có thể tạo mã tạm thời gồm bảy chữ số. Ứng dụng gửi bản tin MQTT qua OneIoT khi có kết nối trong phiên; gửi bản tin chưa chứng minh khóa thật đã áp dụng mã. Backend, các API nghiệp vụ, cơ sở dữ liệu H2 và đường app–server đã có trong mã nguồn; việc xác nhận khả năng đồng bộ trên hạ tầng triển khai liên tục và nhiều thiết bị vẫn cần kiểm thử đầu cuối. Bản Release cũng đã được kiểm tra chạy ở chế độ local.
\end{document}
